\chapter{Modern PHP 8.1--8.4 Engineering: Runtime Internals, Type Systems, and Resilient Architecture}
\label{chap:modern-php}

\begin{chaptergoalsbox}
\textbf{What You Will Learn:}
\begin{itemize}[leftmargin=12pt,itemsep=2pt]
  \item The execution mechanics of PHP-FPM, FastCGI binary records, process pooling, and the shared-nothing lifecycle.
  \item Zend Engine 4 internals: lexical scanning, AST compilation, bytecode emission, and VM opcode dispatch loops.
  \item OPcache shared memory architecture, class preloading, and Tracing JIT machine code generation via DynASM.
  \item Modern type theory in PHP 8.1--8.4: strict types, Union, Intersection, DNF types, and asymmetric property visibility.
  \item Memory management internals: the 16-byte \texttt{zval} struct layout, reference counting, Copy-on-Write (CoW), and the cyclic garbage collection purple-coloring algorithm.
  \item First-class callables, closure binding scopes via \texttt{Closure::bind}, and PSR-7/PSR-15 middleware architecture.
  \item Production webhook systems: constant-time HMAC-SHA256 verification, replay window mitigation, distributed Redis idempotency, and jittered exponential backoff.
\end{itemize}

\textbf{What You Should Already Understand:}
\begin{itemize}[leftmargin=12pt,itemsep=2pt]
  \item Foundational C syntax and memory pointers (\texttt{malloc}, \texttt{struct}, \texttt{union}).
  \item Basic client-server HTTP architectures and web server reverse-proxy mechanics (Nginx/Caddy to application backends).
\end{itemize}

\textbf{Why This Topic Exists:}
Contrary to legacy perceptions, modern PHP (versions 8.1 through 8.4) represents a high-performance, strictly typed systems platform powering foundational web infrastructure (including Wikipedia, Slack backends, Meta Hack lineage, and Stripe-scale payment integrations). Mastering modern PHP requires tearing through high-level abstractions to understand how the Zend Virtual Machine allocates memory, executes bytecode, optimizes hot traces, and achieves zero-shared-state process isolation.

\textbf{Motivating Practical Problem:}
A financial services platform processes incoming payment webhooks from Stripe and Adyen at 4,500 requests per second. Under burst loads, worker processes experience intermittent memory exhaustion, subtle data leakage across concurrent requests, and duplicate transaction execution caused by network replays and missing idempotency locks. Solving this requires understanding the FPM process model, tuning memory allocations down to the \texttt{zval} level, and implementing cryptographically bulletproof, idempotent request pipelines.
\end{chaptergoalsbox}

\section{The PHP-FPM Execution Model and FastCGI Protocol}

At the core of production PHP deployments is **PHP-FPM (FastCGI Process Manager)**. Unlike application runtimes that host long-lived in-memory HTTP servers within a single process (such as Node.js or Python asyncio), standard PHP operates on a **shared-nothing architecture**. Every HTTP request is isolated to an independent worker process, executed from a clean heap state, and torn down upon completion.

\subsection{FastCGI Binary Protocol Framing}

Communication between the reverse proxy (such as Nginx, Caddy, or Envoy) and PHP-FPM occurs over the FastCGI binary protocol (typically across UNIX domain sockets or loopback TCP sockets). FastCGI multiplexes connections using fixed 8-byte record headers:

\begin{lstlisting}[language=C,caption={FastCGI Protocol Header Struct}]
typedef struct {
    unsigned char version;          /* FCGI_VERSION_1 (0x01) */
    unsigned char type;             /* e.g., FCGI_BEGIN_REQUEST (1), FCGI_PARAMS (4), FCGI_STDIN (5) */
    unsigned char requestIdB1;      /* Request ID high byte */
    unsigned char requestIdB0;      /* Request ID low byte */
    unsigned char contentLengthB1;  /* Payload length high byte */
    unsigned char contentLengthB0;  /* Payload length low byte */
    unsigned char paddingLength;    /* Padding byte count */
    unsigned char reserved;
} fcgi_header;
\end{lstlisting}

When an HTTP request arrives, Nginx translates HTTP headers into a stream of \texttt{FCGI\_PARAMS} key-value records, streams the request body inside \texttt{FCGI\_STDIN} records, and terminates the stream with an empty \texttt{FCGI\_STDIN} record. The PHP-FPM worker decodes these records into the superglobals (\texttt{\$\_SERVER}, \texttt{\$\_POST}, \texttt{php://input}).

\subsection{Process Supervision: Master and Worker Pools}

The PHP-FPM process hierarchy consists of a single privileged **Master Process** and multiple unprivileged **Worker Processes** grouped into pools:
\begin{enumerate}[leftmargin=15pt]
  \item \textbf{Master Process:} Runs as \texttt{root} (or a dedicated system user), binds to the listening socket, parses configuration, handles OS signals (\texttt{SIGQUIT} for graceful stop, \texttt{SIGUSR1} for log reopening, \texttt{SIGUSR2} for graceful reload), and supervises workers using \texttt{waitpid()}.
  \item \textbf{Worker Pool:} Workers execute as unprivileged users (e.g., \texttt{www-data}). In dynamic pooling mode, workers scale according to traffic:
  \begin{itemize}
    \item \texttt{pm = dynamic}: FPM maintains \texttt{pm.min\_spare\_servers} and scales up to \texttt{pm.max\_children}.
    \item \texttt{pm = static}: FPM pre-forks a fixed number of workers (\texttt{pm.max\_children}). In high-throughput production, static allocation avoids fork overhead and page-table churn.
  \end{itemize}
\end{enumerate}

\begin{figure}[htbp]
\centering
\begin{tikzpicture}[
  scale=0.95, transform shape,
  node distance=1.2cm and 1.8cm,
  box/.style={rectangle, draw=headingblue, fill=codebg, rounded corners=2pt, thick, align=center, inner sep=7pt, font=\small\sffamily},
  accentbox/.style={rectangle, draw=intuitionTeal, fill=intuitionBg, rounded corners=2pt, thick, align=center, inner sep=7pt, font=\small\sffamily\bfseries},
  line/.style={draw=headingblue, -{Stealth[length=2.5mm]}, thick}
]
  \node[box] (nginx) {\textbf{Reverse Proxy}\\\small Nginx / Envoy};
  \node[accentbox, right=of nginx] (master) {\textbf{FPM Master Process}\\\footnotesize PID 1001 (Binds Socket)};
  \node[box, below=1.0cm of master, xshift=-2.6cm] (w1) {\textbf{Worker 1}\\\footnotesize Request Lifecyle};
  \node[box, below=1.0cm of master] (w2) {\textbf{Worker 2}\\\footnotesize Idle / Accept};
  \node[box, below=1.0cm of master, xshift=2.6cm] (w3) {\textbf{Worker $N$}\\\footnotesize Active Execution};

  \draw[line] (nginx) -- node[above, font=\footnotesize\ttfamily] {FastCGI Stream} (master);
  \draw[line] (master) -- node[left, font=\tiny\sffamily] {fork / epoll} (w1);
  \draw[line] (master) -- (w2);
  \draw[line] (master) -- node[right, font=\tiny\sffamily] {supervise} (w3);

  \node[draw=gray!50, dashed, fit=(w1) (w2) (w3), inner sep=10pt, label=below:{\footnotesize\sffamily Isolated Process Worker Pool (Shared-Nothing Memory)}] {};
\end{tikzpicture}
\caption{Architecture of the PHP-FPM process manager and worker pool isolation.}
\label{fig:fpm-arch}
\end{figure}

\subsection{The Request Lifecycle: Zero-Shared-State Mechanics}

The lifecycle within each worker proceeds through five discrete stages:
\begin{enumerate}[leftmargin=15pt]
  \item \textbf{Module Init (MINIT):} Invoked once when the worker forks. C extensions initialize persistent structures and register functions.
  \item \textbf{Request Init (RINIT):} Executed prior to processing a request. The Zend Memory Manager (ZMM) allocates a request heap, resets global state, initializes superglobals, and resets tick counters.
  \item \textbf{Script Execution:} The entry script (e.g., \texttt{index.php}) is compiled or fetched from OPcache, and executed on the virtual machine.
  \item \textbf{Request Shutdown (RSHUTDOWN):} User output buffers are flushed, registered shutdown functions (\texttt{register\_shutdown\_function}) execute, object destructors are invoked, and the ZMM frees the entire request heap in one bulk operation.
  \item \textbf{Module Shutdown (MSHUTDOWN):} Triggered when the worker process terminates.
\end{enumerate}

\begin{underthehoodbox}[The Immense Engineering Benefit of Bulk Heap Teardown]
In languages with shared persistent heaps (such as Node.js, Python, or Go), a dangling pointer, unclosed connection, or cyclic data structure permanently leaks RAM into the process heap until restart. In PHP-FPM, when \texttt{RSHUTDOWN} completes, the Zend Memory Manager resets its per-request memory arenas and returns free chunks to its internal slab cache, guaranteeing that single-request leaks cannot degrade subsequent executions.
\end{underthehoodbox}

\begin{productionwarningbox}[Process Recycled via \texttt{pm.max\_requests}]
While request heaps are torn down after every request, persistent C extensions or internal OPcache fragmentation can cause slow worker memory drift. In production, always configure:
\begin{lstlisting}[language=sh]
pm.max_requests = 10000
\end{lstlisting}
This forces the FPM master to cleanly terminate and respawn a worker after processing 10,000 requests, preventing memory fragmentation from degrading server throughput.
\end{productionwarningbox}

% ==========================================
\section{Zend Engine Architecture: Lexing to Virtual Machine}

The **Zend Engine 4** is the execution kernel of PHP. It transforms high-level source text into machine-executable virtual instructions through a four-stage pipeline:

\begin{figure}[htbp]
\centering
\begin{tikzpicture}[
  scale=0.9, transform shape,
  node distance=1.0cm and 1.2cm,
  stage/.style={rectangle, draw=headingblue, fill=codebg, rounded corners=2pt, thick, align=center, inner sep=6pt, font=\small\sffamily},
  astage/.style={rectangle, draw=intuitionTeal, fill=intuitionBg, rounded corners=2pt, thick, align=center, inner sep=6pt, font=\small\sffamily\bfseries},
  line/.style={draw=headingblue, -{Stealth[length=2.5mm]}, thick}
]
  \node[stage] (src) {\textbf{Source Code}\\\texttt{.php script}};
  \node[stage, right=of src] (lex) {\textbf{Lexer (re2c)}\\\small Token Stream};
  \node[stage, right=of lex] (parse) {\textbf{Parser (Bison)}\\\small Abstract Syntax Tree};
  \node[astage, below=1.0cm of parse] (comp) {\textbf{Compiler}\\\small Opcode Generation};
  \node[stage, left=of comp] (opcache) {\textbf{OPcache Optimizer}\\\small SSA Pass \& Dead Code};
  \node[astage, left=of opcache] (vm) {\textbf{Zend VM Executor}\\\small Opcode Execution Loop};

  \draw[line] (src) -- (lex);
  \draw[line] (lex) -- (parse);
  \draw[line] (parse) -- (comp);
  \draw[line] (comp) -- (opcache);
  \draw[line] (opcache) -- (vm);
\end{tikzpicture}
\caption{The Zend Engine compilation and execution pipeline.}
\label{fig:zend-pipeline}
\end{figure}

\subsection{Lexical Scanning and AST Construction}
\begin{itemize}[leftmargin=15pt]
  \item \textbf{Lexing:} Implemented via \texttt{Zend/zend\_language\_scanner.l} (compiled using \texttt{re2c}). It converts raw UTF-8 streams into discrete tokens (e.g., \texttt{T\_CLASS}, \texttt{T\_VARIABLE}, \texttt{T\_OBJECT\_OPERATOR}).
  \item \textbf{Parsing:} Implemented in \texttt{Zend/zend\_language\_parser.y} (compiled with GNU Bison LALR(1)). It validates grammar rules and generates the Abstract Syntax Tree (AST), represented by \texttt{zend\_ast} nodes.
\end{itemize}

\subsection{Opcode Compilation and the \texttt{zend\_op} Structure}

The Zend compiler traverses the AST and produces linear arrays of opcodes contained within a \texttt{zend\_op\_array}. A single opcode instruction is represented by the C struct \texttt{zend\_op}:

\begin{lstlisting}[language=C,caption={The \texttt{zend\_op} Internal Representation}]
struct _zend_op {
    const void *handler;     /* Function pointer to VM opcode execution routine */
    znode_op op1;            /* First operand */
    znode_op op2;            /* Second operand */
    znode_op result;         /* Result destination operand */
    uint32_t extended_value; /* Auxiliary instruction flag / operand info */
    uint32_t lineno;         /* Source line number for debugging and stack traces */
    zend_uchar opcode;       /* Opcode number (e.g., ZEND_ADD, ZEND_FETCH_OBJ_R) */
    zend_uchar op1_type;     /* IS_CONST, IS_TMP_VAR, IS_VAR, IS_CV, or IS_UNUSED */
    zend_uchar op2_type;
    zend_uchar result_type;
};
\end{lstlisting}

\subsubsection{Operand Storage Types}
Operands in the Zend Virtual Machine are categorized into five distinct types:
\begin{enumerate}[leftmargin=15pt]
  \item \texttt{IS\_CONST}: Literal constant compiled into the script (e.g., string literals, integer numbers).
  \item \texttt{IS\_TMP\_VAR}: Ephemeral intermediate calculation variable (allocated on the VM evaluation frame, no refcounting needed).
  \item \texttt{IS\_VAR}: Complex VM variable representing objects, references, or property fetches.
  \item \texttt{IS\_CV} (Compiled Variable): Fast local variable (e.g., \texttt{\$total}, \texttt{\$user}). CVs are cached directly in an indexed array on the execution stack, bypassing hash-table lookups entirely!
  \item \texttt{IS\_UNUSED}: Operand slot unused by the specific opcode instruction.
\end{enumerate}

\subsection{The Virtual Machine Dispatch Loop}

The Zend VM executes instructions sequentially using a configurable dispatch mechanism:
\begin{itemize}[leftmargin=15pt]
  \item \textbf{Call Dispatch:} A standard C \texttt{switch(op->opcode)} or function pointer table. Simple, but suffers from high branch misprediction penalties.
  \item \textbf{Direct Threading (GCC/Clang):} The VM uses GNU C \texttt{\&\&label} computed \texttt{goto} arrays. Each opcode handler ends with an indirect jump directly to the next handler address:
  \begin{lstlisting}[language=C]
#define ZEND_VM_NEXT_OPCODE() \
    CHECK_SYMBOL_TABLES() \
    execute_data->opline++; \
    ZEND_VM_CONTINUE()
  \end{lstlisting}
  Direct threading eliminates the central loop dispatcher, drastically reducing CPU pipeline stalls.
\end{itemize}

% ==========================================
\section{OPcache and the JIT (Just-In-Time) Compiler}

Parsing source code and compiling ASTs into opcodes on every HTTP request creates massive CPU overhead. **OPcache** eliminates this by caching compiled \texttt{zend\_op\_array} structures directly inside shared memory (SHM).

\subsection{OPcache Shared Memory Architecture and Preloading}

When a PHP file is requested:
\begin{enumerate}[leftmargin=15pt]
  \item OPcache computes the hash of the real canonical file path.
  \item It checks shared memory for an existing compiled \texttt{zend\_op\_array}.
  \item If present and valid (checked via timestamp comparison or disabled file validation in production), execution begins immediately without touching the disk or lexer.
\end{enumerate}

\subsubsection{PHP Preloading (\texttt{opcache.preload})}
Introduced in PHP 7.4 and enhanced in 8.1+, preloading compiles a specified set of classes, interfaces, and functions into memory at FPM server startup, permanently baking them into process memory before any request arrives:

\begin{lstlisting}[language=PHP,caption={Production Preload Script (\texttt{preload.php})}]
<?php
declare(strict_types=1);

// Recursively traverse and compile core vendor & domain files
$files = new RecursiveIteratorIterator(
    new RecursiveDirectoryIterator(__DIR__ . '/src')
);

foreach ($files as $file) {
    if ($file->isFile() && $file->getExtension() === 'php') {
        opcache_compile_file($file->getRealPath());
    }
}
\end{lstlisting}

Preloaded classes become globally accessible to all workers without autoloader overhead, slashing memory allocation per request.

\subsection{The Tracing JIT Architecture}

PHP 8.0 introduced the **Just-In-Time (JIT)** compiler, developed using **DynASM**. Unlike a method-based JIT that compiles entire functions upon invocation, PHP 8.1+ implements a **Tracing JIT**:

\begin{figure}[htbp]
\centering
\begin{tikzpicture}[
  scale=0.9, transform shape,
  node distance=1.1cm and 1.5cm,
  node/.style={rectangle, draw=headingblue, fill=codebg, rounded corners=2pt, thick, align=center, inner sep=6pt, font=\small\sffamily},
  accent/.style={rectangle, draw=warningRed, fill=warningBg, rounded corners=2pt, thick, align=center, inner sep=6pt, font=\small\sffamily\bfseries},
  line/.style={draw=headingblue, -{Stealth[length=2.5mm]}, thick}
]
  \node[node] (vm) {\textbf{Zend VM Opcode Execution}};
  \node[node, right=of vm] (counter) {\textbf{Loop / Call Profile Counters}};
  \node[accent, right=of counter] (trace) {\textbf{Hot Trace Recorder}\\\small Record sequence of executed opcodes};
  \node[node, below=1.2cm of trace] (ir) {\textbf{Linear Intermediate Representation (IR)}\\\small Type Specialization \& Range Analysis};
  \node[accent, left=of ir] (dynasm) {\textbf{DynASM Machine Code Generation}\\\small Native x86-64 / ARM64 Assembler};
  \node[node, left=of dynasm] (cache) {\textbf{Native Code Cache}\\\small Direct CPU Execution Bypass};

  \draw[line] (vm) -- (counter);
  \draw[line] (counter) -- node[above, font=\tiny\sffamily] {Threshold Exceeded} (trace);
  \draw[line] (trace) -- (ir);
  \draw[line] (ir) -- (dynasm);
  \draw[line] (dynasm) -- (cache);
  \draw[line] (cache) -- (vm);
\end{tikzpicture}
\caption{PHP Tracing JIT compilation flow: identifying hot paths and emitting native machine code.}
\label{fig:jit-flow}
\end{figure}

\subsubsection{JIT Buffer Configuration}
The JIT is configured via the four-digit hexadecimal integer \texttt{opcache.jit}:
\begin{lstlisting}[language=sh]
; CRTO: CPU-specific optimization (1), Register allocation (2), Trigger (0=profile, 5=tracing), Optimization level (4)
opcache.jit = 1255
opcache.jit_buffer_size = 128M
\end{lstlisting}

\begin{intuitionbox}[When Does the JIT Actually Help?]
Traditional web workloads spend over 80\% of their execution time waiting on database queries, network sockets, serialization, and disk I/O. For standard web apps, the JIT yields modest improvements (3--8\%). 

However, for **CPU-bound operations**—such as cryptographic hashing, 3D math, image manipulation, vector search calculations, and complex recursive business algorithms—the Tracing JIT generates raw machine instructions that execute at near-C speeds.
\end{intuitionbox}

% ==========================================
\section{Type System Evolution: PHP 8.1 through 8.4}

PHP's type engine has evolved from a dynamically coerced, weak typing runtime into an expressive, compile-time and runtime verified static-analytical type system.

\subsection{Strict Types and Scalar Type Coercion}

By default, PHP coerces parameter types (e.g., passing string \texttt{"42"} to an \texttt{int} parameter succeeds). Enforcing \texttt{declare(strict\_types=1);} locks strict type checking for the file scope:

\begin{lstlisting}[language=PHP,caption={Enforcing Strict Typing}]
<?php
declare(strict_types=1);

function computeInterest(int $principal, float $rate): float {
    return $principal * $rate;
}

// computeInterest("1000", 0.05); // Fatal Error: TypeError
\end{lstlisting}

\subsection{Union, Intersection, and Disjunctive Normal Form (DNF) Types}

\begin{itemize}[leftmargin=15pt]
  \item \textbf{Union Types (PHP 8.0):} Represents a value that satisfies at least one type: \texttt{int|float|string}.
  \item \textbf{Intersection Types (PHP 8.1):} Requires a value to satisfy all listed types simultaneously (interfaces and classes only): \texttt{Traversable\&Countable}.
  \item \textbf{DNF Types (PHP 8.2):} Combines Union and Intersection types by enforcing Disjunctive Normal Form (an OR of ANDs): \texttt{(A\&B)|C}:
\end{itemize}

\begin{lstlisting}[language=PHP,caption={DNF Types in Domain Repositories}]
interface Entity { public function getId(): string; }
interface Auditable { public function getAuditLog(): array; }
interface NullObject {}

class Repository {
    // Parameter must be (Entity AND Auditable) OR a NullObject
    public function persist((Entity&Auditable)|NullObject $target): void {
        if ($target instanceof NullObject) {
            return;
        }
        // Both Entity and Auditable methods are guaranteed safe by runtime
        $id = $target->getId();
        $log = $target->getAuditLog();
    }
}
\end{lstlisting}

\subsection{Readonly Classes and Enums}

PHP 8.1 introduced pure and backed Enums, which are full objects supporting interfaces and methods. PHP 8.2 introduced \texttt{readonly} classes, enforcing absolute immutability:

\begin{lstlisting}[language=PHP,caption={Backed Enums with Pattern Matching and Readonly Value Objects}]
<?php
declare(strict_types=1);

enum TransactionStatus: string {
    case Pending = 'PENDING';
    case Settled = 'SETTLED';
    case Reversed = 'REVERSED';

    public function isFinal(): bool {
        return match($this) {
            self::Settled, self::Reversed => true,
            self::Pending => false,
        };
    }
}

readonly class Money {
    public function __construct(
        public int $amountInCents,
        public string $currency
    ) {}

    public function add(Money $other): self {
        if ($this->currency !== $other->currency) {
            throw new InvalidArgumentException("Currency mismatch");
        }
        return new self($this->amountInCents + $other->amountInCents, $this->currency);
    }
}
\end{lstlisting}

\subsection{Asymmetric Visibility in PHP 8.4}

PHP 8.4 introduces asymmetric visibility, allowing property reads to be \texttt{public} while restricting property mutations to \texttt{private} or \texttt{protected}:

\begin{lstlisting}[language=PHP,caption={Asymmetric Visibility in PHP 8.4}]
<?php
declare(strict_types=1);

class BankAccount {
    // Read publicly from anywhere, but writable only inside BankAccount
    public private(set) int $balanceInCents;
    public private(set) TransactionStatus $status;

    public function __construct(int $initialDeposit) {
        $this->balanceInCents = $initialDeposit;
        $this->status = TransactionStatus::Pending;
    }

    public function deposit(int $amount): void {
        if ($amount <= 0) throw new InvalidArgumentException();
        $this->balanceInCents += $amount;
    }
}

$acc = new BankAccount(5000);
echo $acc->balanceInCents; // OK: Reads 5000
// $acc->balanceInCents = 10000; // Compile Error: Cannot modify private(set) property
\end{lstlisting}

% ==========================================
\section{Memory Management: Zvals, Reference Counting, and Copy-on-Write}

Understanding how the Zend Engine represents variables in memory is paramount for diagnosing performance bottlenecks and memory leaks in production.

\subsection{The 16-Byte \texttt{zval} Struct}

In PHP 5, a variable required a 32-byte heap-allocated struct plus a 16-byte memory allocator header, totaling 48 bytes per variable. In PHP 7 and 8, the \texttt{zval} was completely redesigned into an ultra-compact **16-byte struct** allocated directly on the CPU stack or inline within hash-table buckets:

\begin{lstlisting}[language=C,caption={Zend \texttt{zval} Struct Definition (\texttt{Zend/zend\_types.h})}]
struct _zval_struct {
    zend_value value;            /* 8-byte union storing raw data or pointer */
    union {
        struct {
            ZEND_ENDIAN_LOHI_3(
                zend_uchar type,         /* Variable type tag (e.g. IS_STRING, IS_ARRAY, IS_OBJECT) */
                zend_uchar type_flags,   /* Type flags (e.g. IS_TYPE_REFCOUNTED, IS_TYPE_COLLECTABLE) */
                union {
                    uint16_t extra;      /* Auxiliary flags */
                } u;
            )
        } v;
        uint32_t type_info;      /* Direct 32-bit access to type and flags */
    } u1;
    union {
        uint32_t next;           /* Collision chain in HashTable */
        uint32_t cache_slot;     /* Opcode cache slot */
        uint32_t opline_num;     /* Execution line */
        uint32_t lineno;         /* Line number */
        uint32_t num_args;       /* Number of function arguments */
    } u2;                        /* 4 bytes of auxiliary metadata */
};
\end{lstlisting}

The \texttt{zend\_value} union is defined as:
\begin{lstlisting}[language=C]
typedef union _zend_value {
    zend_long         lval;    /* 64-bit signed integer */
    double            dval;    /* 64-bit IEEE-754 floating point */
    zend_refcounted  *counted; /* Pointer to reference counted payload */
    zend_string      *str;     /* Pointer to string structure */
    zend_array       *arr;     /* Pointer to array HashTable */
    zend_object      *obj;     /* Pointer to object structure */
    zend_resource    *res;     /* Pointer to resource */
    zend_reference   *ref;     /* Pointer to explicit reference */
    zend_ast_ref     *ast;     /* Pointer to AST */
} zend_value;
\end{lstlisting}

\begin{figure}[htbp]
\centering
\begin{tikzpicture}[
  scale=0.95, transform shape,
  node distance=0.2cm,
  bytebox/.style={rectangle, draw=headingblue, fill=codebg, thick, minimum height=1.0cm, align=center, font=\footnotesize\ttfamily},
  headerbox/.style={rectangle, draw=none, align=center, font=\small\sffamily\bfseries}
]
  \node[bytebox, minimum width=6.0cm] (val) {\textbf{value (8 bytes)}\\Stores \texttt{lval}, \texttt{dval}, or \texttt{*counted} pointer};
  \node[bytebox, minimum width=3.0cm, right=of val] (u1) {\textbf{u1 (4 bytes)}\\\texttt{type}, \texttt{type\_flags}};
  \node[bytebox, minimum width=3.0cm, right=of u1] (u2) {\textbf{u2 (4 bytes)}\\\texttt{cache\_slot} / \texttt{next}};

  \node[draw=intuitionTeal, thick, dashed, fit=(val) (u1) (u2), inner sep=4pt, label=above:{\textbf{\sffamily Exact 16-Byte Inline \texttt{zval} Layout}}] {};
\end{tikzpicture}
\caption{The 16-byte \texttt{zval} memory layout in PHP 8.}
\label{fig:zval-layout}
\end{figure}

Notice that primitive scalars (\texttt{int}, \texttt{float}, \texttt{bool}, \texttt{null}) require **zero external heap allocations**! Their values live entirely inside the 8-byte \texttt{zend\_value} union.

\subsection{Reference Counting Mechanics}

Complex structures (\texttt{string}, \texttt{array}, \texttt{object}, \texttt{resource}) cannot fit within 8 bytes. Instead, the \texttt{zval} holds a pointer to a heap-allocated header conforming to \texttt{zend\_refcounted}:

\begin{lstlisting}[language=C]
typedef struct _zend_refcounted_h {
    uint32_t         refcount; /* Number of references pointing to this memory */
    union {
        struct {
            ZEND_ENDIAN_LOHI_3(
                zend_uchar    type,
                zend_uchar    flags,
                uint16_t      gc_info  /* Garbage collection color & root buffer pointer */
            )
        } v;
        uint32_t type_info;
    } u;
} zend_refcounted_h;
\end{lstlisting}

Every time a variable is assigned or passed by value to a function, the engine does not clone the data. It simply increments \texttt{refcount++}.

\subsection{Copy-on-Write (CoW) Mechanics}

Copy-on-Write guarantees maximum memory efficiency by deferring duplication until the moment of mutation:

\begin{lstlisting}[language=PHP,caption={Demonstrating Copy-on-Write Behavior}]
<?php
$a = range(1, 100_000); // Allocates ~4MB array. $a->value->arr->refcount = 1
$b = $a;                // Zero extra memory allocated! $a and $b point to SAME arr.
                        // $a->value->arr->refcount = 2

$b[] = 100_001;         // Mutation! CoW triggers.
                        // The engine checks: is refcount > 1?
                        // YES -> Clones the array!
                        // $b points to new arr (refcount = 1).
                        // $a points to original arr (refcount = 1).
\end{lstlisting}

\subsection{Interned and Immutable Strings}

Strings known at compile time (such as function names, class names, and constant string literals) are allocated as **interned strings**. Their \texttt{IS\_STR\_INTERNED} flag is set, and their \texttt{refcount} is initialized to a static constant that is never incremented or decremented. They reside in shared memory and are never freed during request teardown.

% ==========================================
\section{Circular Garbage Collection: The Roots Buffer and Coloring Algorithm}

Reference counting has one fatal mathematical limitation: **cyclic references**. If object $A$ holds a reference to object $B$, and object $B$ holds a reference to object $A$, their \texttt{refcount} values will never drop to zero, even when all external references are unset.

\subsection{The Cyclic Garbage Collector Algorithm}

To solve this, PHP implements a cyclic garbage collector based on the **Bacon and Rajan algorithm**:

\begin{figure}[htbp]
\centering
\begin{tikzpicture}[
  scale=0.9, transform shape,
  node distance=1.0cm and 1.2cm,
  state/.style={circle, draw=headingblue, thick, minimum size=1.8cm, align=center, font=\small\sffamily},
  line/.style={draw=headingblue, -{Stealth[length=2.5mm]}, thick}
]
  \node[state, fill=purple!20] (purple) {\textbf{Purple}\\\footnotesize Possible Root};
  \node[state, fill=gray!30, right=of purple] (gray) {\textbf{Grey}\\\footnotesize Traced Decr};
  \node[state, fill=black, text=white, below=1.2cm of gray] (black) {\textbf{Black}\\\footnotesize Live Object};
  \node[state, fill=white, draw=gray, below=1.2cm of purple] (white) {\textbf{White}\\\footnotesize Garbage Candidate};

  \draw[line] (purple) -- node[above, font=\tiny\sffamily] {gc\_collect\_cycles()} (gray);
  \draw[line] (gray) -- node[right, font=\tiny\sffamily] {refcount == 0} (white);
  \draw[line] (gray) -- node[right, font=\tiny\sffamily] {refcount > 0} (black);
  \draw[line] (white) -- node[left, font=\tiny\sffamily] {zval\_dtor()} (purple);
\end{tikzpicture}
\caption{The Garbage Collector cycle detection coloring states.}
\label{fig:gc-states}
\end{figure}

The collector operates through four deterministic phases:
\begin{enumerate}[leftmargin=15pt]
  \item \textbf{Mark Purple:} When a refcount is decremented but does not reach zero, it might be part of an orphan cycle. The pointer is colored **Purple** and appended to the **Roots Buffer** (capacity: 10,000 entries by default).
  \item \textbf{Mark Grey (Trial Deletion):} When the roots buffer fills, the GC triggers. It performs a depth-first traversal of all children, decrementing their internal refcounts by 1 and coloring visited nodes **Grey**.
  \item \textbf{Scan Roots (Mark Black or White):}
  \begin{itemize}
    \item If a node's refcount is strictly greater than 0, it must be reachable from an external scope. The GC rolls back the decrements and colors the subgraph **Black** (alive).
    \item If a node's refcount reached exactly 0, it belongs entirely to an isolated cyclic loop. It is colored **White** (dead garbage).
  \end{itemize}
  \item \textbf{Sweep (Freeing):} All White nodes are freed from memory and their buffers returned to the allocator.
\end{enumerate}

\begin{productionwarningbox}[Garbage Collection in Long-Running Daemons]
In short-lived web requests, cyclic GC can often be disabled (\texttt{zend.enable\_gc = 0}) because \texttt{RSHUTDOWN} frees the entire heap at request completion. 

However, in **long-running CLI workers** (e.g., Laravel Queue workers, Symfony Messenger, RoadRunner, FrankenPHP), failing to monitor GC buffers will lead to sudden stop-the-world GC pauses or unbounded heap bloat. If memory steadily creeps upward despite \texttt{unset()}, an uncollected cycle is holding the root!
\end{productionwarningbox}

% ==========================================
\section{Closures, Scoping, and First-Class Callables}

Functions in PHP are first-class language constructs. Closures are instantiated as internal \texttt{Closure} class objects.

\subsection{Closure State and Scope Binding}

When a closure is created, it captures lexical scope via the \texttt{use} keyword:
\begin{lstlisting}[language=PHP]
$taxRate = 0.20;
$calculator = function(int $amount) use ($taxRate): float {
    return $amount * (1 + $taxRate);
};
\end{lstlisting}

Internally, \texttt{\$taxRate} is copied and stored within the \texttt{Closure} object's private properties table.

\subsection{Rebinding Scope with \texttt{Closure::bind}}

PHP empowers dynamic metaprogramming by allowing closures to rebind their execution scope and bound object instance at runtime:

\begin{lstlisting}[language=PHP,caption={Extracting Private State via \texttt{Closure::bindTo}}]
<?php
declare(strict_types=1);

class SecretVault {
    private string $encryptionKey = "SUPER_SECRET_AES_KEY_0xFA";
}

$vault = new SecretVault();

// Create an unbound closure
$extractor = function() {
    return $this->encryptionKey; // Accesses private property
};

// Bind closure to the instance and grant scope access to SecretVault class
$boundExtractor = $extractor->bindTo($vault, SecretVault::class);
echo $boundExtractor(); // Prints: SUPER_SECRET_AES_KEY_0xFA
\end{lstlisting}

\subsection{First-Class Callable Syntax (PHP 8.1+)}

In PHP 8.1+, any function, static method, or instance method can be converted into a standard \texttt{Closure} using the syntax \texttt{fn(...)}:

\begin{lstlisting}[language=PHP,caption={First-Class Callable Examples}]
$hasher = hash_hmac(...); // First-class callable
$result = $hasher('sha256', 'payload', 'secret');

$validator = $this->validateTransaction(...);
array_walk($transactions, $validator);
\end{lstlisting}

% ==========================================
\section{PSR Standards: Architectural Foundations}

The PHP Standard Recommendations (PSR), authored by the PHP-FIG (Framework Interoperability Group), define the structural contracts that enable modern ecosystem interoperability.

\subsection{PSR-4 Autoloader Specification}

PSR-4 defines how fully-qualified class names map directly to file-system paths. The formal specification dictates:
$$\text{Class: } \langle\text{NamespacePrefix}\rangle\backslash\langle\text{SubNamespaces}\rangle\backslash\langle\text{ClassName}\rangle$$
Maps directly to:
$$\langle\text{BaseDirectory}\rangle/\langle\text{SubDirectoryPath}\rangle/\langle\text{ClassName}\rangle\text{.php}$$

Composer compiles these mappings into an optimized static lookup table:
\begin{lstlisting}[language=sh]
composer dump-autoload --optimize --classmap-authoritative
\end{lstlisting}
This bypasses disk filesystem \texttt{file\_exists()} calls entirely by evaluating an in-memory hash table in shared memory.

\subsection{PSR-7 and PSR-15 Middleware Pipelines}

\begin{itemize}[leftmargin=15pt]
  \item \textbf{PSR-7 (HTTP Message Interfaces):} Models HTTP requests and responses as **immutable value objects**. Mutating a request returns a cloned instance:
  \begin{lstlisting}[language=PHP]
  $newRequest = $request->withHeader('X-Correlation-ID', $uuid);
  \end{lstlisting}
  \item \textbf{PSR-15 (HTTP Handlers and Middleware):} Defines the standard onion-layer middleware pipeline:
\end{itemize}

\begin{lstlisting}[language=PHP,caption={PSR-15 Middleware Implementation}]
use Psr\Http\Message\ResponseInterface;
use Psr\Http\Message\ServerRequestInterface;
use Psr\Http\Server\MiddlewareInterface;
use Psr\Http\Server\RequestHandlerInterface;

class CorrelationIdMiddleware implements MiddlewareInterface {
    public function process(
        ServerRequestInterface $request,
        RequestHandlerInterface $handler
    ): ResponseInterface {
        $correlationId = $request->getHeaderLine('X-Correlation-ID') ?: bin2hex(random_bytes(16));
        $requestWithId = $request->withHeader('X-Correlation-ID', $correlationId);
        
        $response = $handler->handle($requestWithId);
        
        return $response->withHeader('X-Correlation-ID', $correlationId);
    }
}
\end{lstlisting}

% ==========================================
\section{Production Webhook Engineering: Signatures, Replay Defense, and Idempotency}

Building enterprise webhook consumers (e.g., Stripe, Shopify, GitHub) requires a fault-tolerant architecture capable of handling out-of-order delivery, network replays, and forged payloads.

\begin{figure}[htbp]
\centering
\begin{tikzpicture}[
  scale=0.9, transform shape,
  node distance=1.0cm and 1.2cm,
  box/.style={rectangle, draw=headingblue, fill=codebg, rounded corners=2pt, thick, align=center, inner sep=6pt, font=\small\sffamily},
  accent/.style={rectangle, draw=intuitionTeal, fill=intuitionBg, rounded corners=2pt, thick, align=center, inner sep=6pt, font=\small\sffamily\bfseries},
  line/.style={draw=headingblue, -{Stealth[length=2.5mm]}, thick}
]
  \node[box] (sender) {\textbf{Stripe / Webhook Sender}};
  \node[accent, right=of sender] (verify) {\textbf{HMAC-SHA256 Verification}\\\small Constant-time \texttt{hash\_equals}};
  \node[box, right=of verify] (replay) {\textbf{Timestamp Replay Filter}\\\small Reject if $|t_{\text{now}} - t_{\text{msg}}| > 300\text{s}$};
  \node[accent, below=1.2cm of replay] (idem) {\textbf{Redis Atomic Mutex}\\\small \texttt{SET key uuid NX EX 86400}};
  \node[box, left=of idem] (queue) {\textbf{Durable Queue (RabbitMQ)}\\\small Asynchronous Worker};
  \node[box, left=of queue] (ack) {\textbf{HTTP 200 OK}\\\small Within 500ms};

  \draw[line] (sender) -- node[above, font=\tiny\sffamily] {HTTP POST} (verify);
  \draw[line] (verify) -- (replay);
  \draw[line] (replay) -- (idem);
  \draw[line] (idem) -- (queue);
  \draw[line] (queue) -- (ack);
\end{tikzpicture}
\caption{Production webhook ingestion pipeline featuring signature check, replay window filter, and Redis idempotency locks.}
\label{fig:webhook-pipeline}
\end{figure}

\subsection{Constant-Time HMAC-SHA256 Verification}

Standard string equality comparisons (\texttt{===}) terminate on the first mismatched byte, creating a measurable **timing side-channel** that allows an attacker to deduce the valid signature byte-by-byte. Production verification requires \texttt{hash\_equals()}, which executes in constant time:

\begin{lstlisting}[language=PHP,caption={Constant-Time HMAC Verification and Replay Defense}]
<?php
declare(strict_types=1);

final readonly class WebhookSignatureValidator {
    private const int ALLOWED_CLOCK_DRIFT_SECONDS = 300; // 5 minute replay window

    public function __construct(
        private string $webhookSecret
    ) {}

    public function validate(string $rawPayload, string $sigHeader): void {
        // Stripe style header: t=1614555000,v1=5257a869e7ec...
        $parts = [];
        foreach (explode(',', $sigHeader) as $item) {
            [$key, $value] = explode('=', trim($item), 2) + ['', ''];
            $parts[$key] = $value;
        }

        if (!isset($parts['t'], $parts['v1'])) {
            throw new RuntimeException("Malformed signature header");
        }

        $timestamp = (int)$parts['t'];
        $signature = $parts['v1'];

        // Replay defense: Reject requests older than tolerance window
        if (abs(time() - $timestamp) > self::ALLOWED_CLOCK_DRIFT_SECONDS) {
            throw new RuntimeException("Webhook timestamp outside allowed tolerance window");
        }

        // Compute expected signature: hash_hmac(sha256, timestamp . "." . payload, secret)
        $signedPayload = $timestamp . '.' . $rawPayload;
        $expectedSignature = hash_hmac('sha256', $signedPayload, $this->webhookSecret);

        // Constant-time execution prevents timing attacks
        if (!hash_equals($expectedSignature, $signature)) {
            throw new RuntimeException("Invalid webhook signature verification failed");
        }
    }
}
\end{lstlisting}

\subsection{Distributed Idempotency with Redis Mutexes}

Network transients cause webhook providers to retry delivery, frequently delivering identical events multiple times. To guarantee single-execution semantics, the ingestion pipeline acquires an atomic Redis lock on the event ID:

\begin{lstlisting}[language=PHP,caption={Atomic Redis Idempotency Lock}]
final readonly class WebhookIdempotencyManager {
    public function __construct(
        private Redis $redis
    ) {}

    public function acquire(string $eventId, int $ttlSeconds = 86400): bool {
        $key = 'webhook:idempotency:' . $eventId;
        
        // SET key val NX EX ttl -> Atomic check-and-set in Redis
        $acquired = $this->redis->set($key, 'PROCESSING', ['NX', 'EX' => $ttlSeconds]);
        
        return (bool)$acquired;
    }

    public function markCompleted(string $eventId): void {
        $key = 'webhook:idempotency:' . $eventId;
        $this->redis->set($key, 'COMPLETED', ['XX', 'KEEPTTL']);
    }
}
\end{lstlisting}

\subsection{Exponential Backoff with Full Jitter}

When dispatching outbound webhooks to third-party endpoints, retrying on a fixed schedule causes **thundering herd** problems that overwhelm recovering servers. Production dispatchers implement exponential backoff with full jitter (randomized sleep):

$$T_{\text{sleep}} = \text{random\_int}(0, \; \min(T_{\text{max}}, \; T_{\text{base}} \times 2^{\text{attempt}}))$$

\begin{lstlisting}[language=PHP,caption={Exponential Backoff with Full Jitter}]
function calculateJitteredBackoff(int $attempt, int $baseMs = 100, int $maxMs = 10000): int {
    $exponentialLimit = min($maxMs, $baseMs * (2 ** $attempt));
    return random_int(0, $exponentialLimit);
}
\end{lstlisting}

% ==========================================
\section{Chapter Summary \& Key Takeaways}

\begin{itemize}[leftmargin=15pt]
  \item \textbf{PHP-FPM Architecture:} Isolates execution per request using a shared-nothing memory lifecycle. The Zend Memory Manager destroys the request heap during \texttt{RSHUTDOWN}, eliminating cross-request heap leakage.
  \item \textbf{Zend VM Internals:} The engine parses code into an AST and compiles it into \texttt{zend\_op\_array} structures. Compiled variables (\texttt{IS\_CV}) reside directly on the execution stack for instantaneous indexed array lookups.
  \item \textbf{OPcache \& JIT:} OPcache eliminates parsing overhead by persisting opcodes in shared memory. The Tracing JIT detects hot loops and emits native x86-64/ARM64 machine instructions via DynASM.
  \item \textbf{Memory Representation:} The 16-byte \texttt{zval} stores scalars inline. Reference counting and Copy-on-Write prevent memory cloning until mutation occurs.
  \item \textbf{Garbage Collection:} The cyclic roots buffer employs a four-color trial deletion algorithm to identify and collect orphaned circular references.
  \item \textbf{Modern Type Safety:} PHP 8.1--8.4 provides strict types, Union, Intersection, DNF types, readonly classes, and asymmetric visibility (\texttt{public private(set)}).
  \item \textbf{Webhook Hardening:} Production webhooks demand constant-time \texttt{hash\_equals()} HMAC validation, clock drift validation, Redis-backed atomic idempotency, and jittered backoff.
\end{itemize}

% ==========================================
\section{Exercises and Real-World Scenarios}

\subsection{Mechanics Exercises}
\begin{enumerate}
  \item \textbf{Exercise 1.1:} Draw the memory state and trace the exact \texttt{refcount} value of the underlying array during each line of the following code snippet:
\begin{lstlisting}[language=PHP]
$x = [1, 2, 3];
$y = $x;
$z = &$y;
$z[] = 4;
\end{lstlisting}
  \item \textbf{Exercise 1.2:} Write a valid PHP 8.4 class \texttt{UserProfile} that utilizes asymmetric property visibility to expose \texttt{\$username} and \texttt{\$email} as publicly readable, but mutable only from within the class or its subclasses, while ensuring they cannot be uninitialized.
\end{enumerate}

\subsection{Debugging Exercises}
\begin{enumerate}
  \item \textbf{Exercise 1.3:} A CLI daemon importing millions of records crashes after 45 minutes with \texttt{Fatal Error: Allowed memory size exhausted}. The code is structured as follows:
\begin{lstlisting}[language=PHP]
while ($row = $db->fetchNext()) {
    $node = new Node($row);
    $node->parent = $currentNode;
    $currentNode->addChild($node);
    $currentNode = $node;
    unset($node);
}
\end{lstlisting}
Explain why \texttt{unset(\$node)} failed to release memory, and provide the exact architectural fix.
\end{enumerate}

\subsection{Architecture Design Problem}
\begin{enumerate}
  \item \textbf{Exercise 1.4:} Architect an end-to-end webhook ingestion gateway in PHP 8.3 that receives payments from Stripe, validates the payload in constant time, prevents replays using a 5-minute sliding window, ensures exactly-once execution using Redis, and enqueues the payload onto a RabbitMQ cluster. Provide the complete implementation.
\end{enumerate}

% ==========================================
\section{Detailed Solutions}

\subsection{Solution 1.1: Tracing Refcounts and References}
\begin{enumerate}[leftmargin=15pt]
  \item \texttt{\$x = [1, 2, 3];}: A \texttt{zend\_array} is allocated on the heap. \texttt{\$x} is a \texttt{zval} of type \texttt{IS\_ARRAY} pointing to the array. \texttt{arr->refcount = 1}.
  \item \texttt{\$y = \$x;}: Copy-on-Write assignment. \texttt{\$y} points to the same \texttt{zend\_array}. \texttt{arr->refcount = 2}.
  \item \texttt{\$z = \&\$y;}: An explicit reference (\texttt{\&}) is created! The engine separates \texttt{\$y} from the simple array pointer. It creates a \texttt{zend\_reference} wrapper with \texttt{refcount = 2} (held by \texttt{\$y} and \texttt{\$z}). The underlying array's refcount is adjusted accordingly.
  \item \texttt{\$z[] = 4;}: Because \texttt{\$x} still points to the original array (with \texttt{refcount > 1}), modifying via reference \texttt{\$z} forces the engine to duplicate (CoW separation) the array so that \texttt{\$x} remains untouched as \texttt{[1, 2, 3]}. \texttt{\$y} and \texttt{\$z} now jointly point to the newly mutated array \texttt{[1, 2, 3, 4]}.
\end{enumerate}

\subsection{Solution 1.2: Asymmetric Visibility Implementation}
\begin{lstlisting}[language=PHP]
<?php
declare(strict_types=1);

class UserProfile {
    public protected(set) string $username;
    public protected(set) string $email;

    public function __construct(string $username, string $email) {
        $this->updateContact($username, $email);
    }

    public function updateContact(string $username, string $email): void {
        if (!filter_var($email, FILTER_VALIDATE_EMAIL)) {
            throw new InvalidArgumentException("Invalid email format");
        }
        $this->username = $username;
        $this->email = $email;
    }
}
\end{lstlisting}

\subsection{Solution 1.3: Debugging Cyclic Memory Leaks}
\textbf{Diagnosis:} Each \texttt{Node} holds a reference to its parent (\texttt{\$node->parent = \$currentNode}), and the parent holds an array containing the child (\texttt{\$currentNode->addChild(\$node)}). This creates a bidirectional cyclic reference graph. 

When \texttt{unset(\$node)} executes, the child's refcount drops from 2 to 1 (because the parent still holds it), and the parent's refcount remains 1. Because the refcount never reaches 0, the Zend Engine cannot free either object. The objects are pushed to the GC Roots Buffer. In high-frequency CLI loops, the cycle generation rate drastically outpaces the collector or causes heavy GC latency spikes.

\textbf{Fix:}
\begin{enumerate}
  \item Use Weak References (\texttt{WeakReference::create(\$currentNode)}) for the child-to-parent pointer, preventing refcount increments.
  \item Or explicitly break the cycle: \texttt{\$node->parent = null;} before unsetting.
  \item Periodically invoke \texttt{gc\_collect\_cycles()} every 5,000 iterations.
\end{enumerate}

\subsection{Solution 1.4: Complete Webhook Ingestion Implementation}
See Section 1.7 for the fully detailed, production-hardened \texttt{WebhookSignatureValidator} and \texttt{WebhookIdempotencyManager} implementations.

% ==========================================
\section{Further Reading \& Citations}

\begin{itemize}[leftmargin=15pt]
  \item Bacon, D. F., \& Rajan, V. T. (2001). \textit{Concurrent Cycle Collection in Reference Counted Systems}. European Conference on Object-Oriented Programming (ECOOP).
  \item Stogov, D. (2020). \textit{The PHP 8 JIT Architecture}. Zend by Perforce Internal Architecture Notes.
  \item PHP-FIG. (2016). \textit{PSR-7: HTTP Message Interfaces} \& \textit{PSR-15: HTTP Server Request Handlers}. \url{https://www.php-fig.org/psr/}
  \item Turnbull, P. (2023). \textit{PHP Internals: Understanding the Zend Engine zval Layout}. PHP Foundation Technical Papers.
\end{itemize}
